% sections/conclusion.tex --- Conclusion
\section{Conclusion}
\label{sec:conclusion}

Empirical machine-learning research on accounting-anomaly detection
has matured to the point where standardised, reproducible
benchmarks are overdue.  The community has accumulated a respectable
toolkit---deep autoencoders, adversarial detectors, gradient
boosting, graph neural networks---but no shared evaluation frame on
which to compare them, because the underlying data is held under
audit-confidentiality constraints that public release does not
respect.  Synthetic data, when it is held to a high empirical and
statistical standard, dissolves that constraint.

\bench{} is an attempt to construct exactly such a shared frame.
It is calibrated against an empirical study of $364$ million
journal-entry postings drawn from $155$ general
ledgers~\citep{ivertowski2024apn}; it is generated by the
open-source \system{} framework~\citep{ivertowski2026datasynth,
datasynth-repo} under a single declarative configuration; it
exposes a $13$-column flat edge list whose schema is fixed and
joinable against the row-level journal-entry table; and it ships
with three reproducible splits, an explicit leakage audit, and a
baseline suite spanning graph neural networks, gradient boosting,
classical detectors, and an ISA~240-style rule-based reference.

\bench{} does not displace empirical work on real client data.
Its role is complementary: a public reference against which methods
can be developed, ablated, compared, and reproduced without the
confidentiality concerns that block public release of the
underlying ledgers.  We anticipate that this role will become more
useful as the benchmark matures: rolling-window splits,
cross-jurisdictional configurations, counterfactual pairs, and
adjacent challenges (control testing, cutoff testing) are all
within reach of the present infrastructure.

We invite the community to use \bench{}, to submit to the
leaderboard, and to file issues against the configuration where
the empirical calibration drifts away from contemporary client
practice.  The benchmark is most valuable when it is held
collectively to the standard the field needs.
